% SGA 3, Expose III, Loop 1 English reconstruction.
% Scope: Lemma 4.17 through Remark 4.22, stopping immediately before 4.23.
% Authority: Polo--Gille Exp3-14oct24.pdf, component-local pp. 63--67
% (combined-reader pp. 173--177; printed pp. 143--148), with only the
% 4.17--4.22 portion of the last page admitted. OCR is a locator only.
% English wording was compared with Jacob C. Reinhold's CC BY 4.0 Markdown
% lineage at commit e7a259f3f8608ad3edf9bf6eead3fd504dd2d23e. The
% Polo--Gille PDF controls prose, mathematics, numbering, editorial notes,
% and the four source-derived Loop-1 diagram images included below.

\setcounter{footnote}{108}

% Source page: Exp3 local p. 63; combined p. 173; printed pp. 143--144.
\medskip
\noindent\SGAInlineRefTarget{sga3:iii:lemma:4-17}\textbf{引理 4.17.}\footnote{编者注: 已修改\SGARefLink{sga3:iii:lemma:4-17}{4.17} 和
\SGARefLink{sga3:iii:paragraph:4-18}{4.18}, 以纳入\SGARefLink{sga3:iii:paragraph:4-16}{4.16} 中增补的内容.}
\textit{设 \(X\) 为在 \(S\) 上平滑的 \(S\)-群,
\(Y_{\mathscr J}\) 为 \(X_{\mathscr J}\) 的子 \(S_{\mathscr J}\)-群,
在 \(S_{\mathscr J}\) 上平坦且局部有限呈示.
对 \(X\) 的每个在 \(S\) 上平坦且提升 \(Y_{\mathscr J}\) 的子概形 \(Y\),
考虑\SGARefLink{sga3:iii:definition:4-16-1}{4.16.1} 中定义的障碍:}
\[
 DY\in
 \operatorname{Hom}_{S_0}
 \bigl(Y_0\mathbin{\times_{S_0}}Y_0,N_0\bigr)
 =C^2(Y_0,N_0).
\]
\begin{enumerate}
\item[\SGAInlineRefTarget{sga3:iii:lemma:4-17:item:i}(i)] \textit{\(Y\) 是 \(X\) 的子 \(S\)-群的充分必要条件为 \(DY=0\).}
\item[\SGAInlineRefTarget{sga3:iii:lemma:4-17:item:ii}(ii)] \textit{若通过函子性, 使 \(Y_0\) 经由 \(Y_0\) 的内自同构作用于 \(N_0\),
则 \(DY\in Z^2(Y_0,N_0)\).}
\item[\SGAInlineRefTarget{sga3:iii:lemma:4-17:item:iii}(iii)] \textit{若 \(Y\) 和 \(Y'\) 为 \(X\) 的两个在 \(S\) 上平坦且提升
\(Y_{\mathscr J}\) 的子概形 (因而偏差
\(d(Y,Y')\in C^1(Y_0,N_0)\) 有定义; 参见\SGARefLink{sga3:iii:remark:4-5-1}{4.5.1}), 则}
  \[
   DY'-DY=\partial^1d(Y,Y').
  \]
  \footnote{编者注: 原版作 \(DY'-DY=-\partial d(Y,Y')\),
但其中的算子 \(\partial\) 是第 I 篇\SGARefLink{sga3:i:subsection:5-1}{5.1}
所定义的微分 \(\partial^1\) 的相反数.}
\end{enumerate}

下面依次证明这些断言.

\medskip
\noindent\SGAInlineRefTarget{sga3:iii:paragraph:4-18}\textbf{4.18. --- 证明 \SGARefLink{sga3:iii:lemma:4-17:item:i}{4.17(i)}.}
由定义, \(DY=0\) 当且仅当 \(Y\) 在 \(X\) 的群运算下稳定.
因此, 若 \(Y\) 为 \(X\) 的子群, 则 \(DY=0\).
反之, 若 \(DY=0\), 则 \(Y\) 带有诱导的运算 \(\mu_Y\),
该运算结合, 且模 \(\mathscr J\) 约化后为 \(X_{\mathscr J}\) 上的群运算.
由于 \(Y\) 在 \(S\) 上平坦且局部有限呈示,
由\SGARefLink{sga3:iii:proposition:3-3}{3.3} 可知 \(\mu_Y\) 为群运算.

\medskip
\noindent\SGAInlineRefTarget{sga3:iii:paragraph:4-19}\textbf{4.19. --- 证明 \SGARefLink{sga3:iii:lemma:4-17:item:ii}{4.17(ii)}.}
比较 \(u=c(X,Y,\mu^2\circ(i,i,i))\)
在下述两个图 \((D_j)\), \(j=1,2\) 中算得的两个值:

% Raster stand-in for Loop 1; replace by a native TeX diagram in Loop 2.
% Authority: Exp3 local p. 63; combined reader p. 173; printed p. 144.
\begin{center}
 \SGAThreeNativeInput[0.94\textwidth]
   {figures/exp3/diagram_056_exp3_p063_Dj_associativity.tex}
\end{center}

这里 \(f_1=(1,\pi)\), \(f_2=(\pi,1)\), 而 \(\mu^2\) 表示态射
\[
 \mu\circ f_1=\mu\circ f_2:
 X\mathbin{\times_S}X\mathbin{\times_S}X\longrightarrow X.
\]
\footnote{编者注: 略微修改了记号, 并扩充了论证的开头.}
令
\[
 \mu_Y=\mu\circ(i,i),\qquad
 f_j^Y=f_j\circ(i,i,i),\qquad
 \mu^{2,Y}=\mu^2\circ(i,i,i).
\]
对 \(j=1,2\), 设 \(a_j\) 和 \(b_j\) 为态射
\[
 a_j=a_{\mu_0}\bigl((f_j^Y)_0\bigr),
 \qquad
 b_j=b_{(f_j^Y)_0},
\]
它们由\SGARefLink{sga3:iii:lemma:4-12}{引理 4.12} 关联于态射对
\((f_j^Y,\mu)\). 因而
\begin{equation}
\tag{\(\dagger\)}\label{sga3:iii:paragraph:4-19:equation:dagger}
\left\{
\begin{aligned}
 a_j:\;&
 \operatorname{Hom}_{\mathscr O_{Y_0^3}}
 \left((f_j^Y)_0^*
  \mathscr N_{Y_0\times Y_0/X_0\times X_0},
  \mathscr J\mathscr O_{Y_0^3}\right)
 \longrightarrow\\[-0.15em]
 &\qquad
 \operatorname{Hom}_{\mathscr O_{Y_0^3}}
 \left((\mu^{2,Y})_0^*\mathscr N_{Y_0/X_0},
  \mathscr J\mathscr O_{Y_0^3}\right),\\[0.35em]
 b_j:\;&
 \operatorname{Hom}_{\mathscr O_{Y_0^2}}
 \left((\mu_Y)_0^*\mathscr N_{Y_0/X_0},
  \mathscr J\mathscr O_{Y_0^2}\right)
 \longrightarrow\\[-0.15em]
 &\qquad
 \operatorname{Hom}_{\mathscr O_{Y_0^3}}
 \left((\mu^{2,Y})_0^*\mathscr N_{Y_0/X_0},
  \mathscr J\mathscr O_{Y_0^3}\right).
\end{aligned}
\right.
\end{equation}

% Source page: Exp3 local p. 64; combined p. 174; printed pp. 144--145.
由于
\[
 \mathscr N_{Y_0\times Y_0/X_0\times X_0}
 \simeq
 \operatorname{pr}_1^*\mathscr N_{Y_0/X_0}
 \oplus
 \operatorname{pr}_2^*\mathscr N_{Y_0/X_0}
\]
(因为 \(X_0\) 和 \(Y_0\) 在 \(S_0\) 上平坦), 且
\[
 \mathscr N_{Y_0/X_0}
 \simeq
 \mathbf n_{Y_0/X_0}\otimes_{\mathscr O_{S_0}}\mathscr O_{Y_0},
\]
有
\[
 (f_j^Y)_0^*\mathscr N_{Y_0\times Y_0/X_0\times X_0}
 \simeq
 (\mathbf n_{Y_0/X_0}\oplus \mathbf n_{Y_0/X_0})
 \otimes\mathscr O_{Y_0^3},
\]
同样有
\[
 \begin{aligned}
 (\mu^{2,Y})_0^*\mathscr N_{Y_0/X_0}
 &\simeq \mathbf n_{Y_0/X_0}\otimes\mathscr O_{Y_0^3},\\
 (\mu_Y)_0^*\mathscr N_{Y_0/X_0}
 &\simeq \mathbf n_{Y_0/X_0}\otimes\mathscr O_{Y_0^2}.
 \end{aligned}
\]
此外, 由于 \(Y_0^2\) 和 \(Y_0^3\) 在 \(S_0\) 上平坦,
\((\dagger)\) 可改写为
\begin{equation}
\tag{\(\ddagger\)}\label{sga3:iii:paragraph:4-19:equation:ddagger}
\left\{
\begin{aligned}
 a_j&:\operatorname{Hom}_{S_0}(Y_0^3,N_0\oplus N_0)
       \longrightarrow\operatorname{Hom}_{S_0}(Y_0^3,N_0),\\
 b_j&:\operatorname{Hom}_{S_0}(Y_0^2,N_0)
       \longrightarrow\operatorname{Hom}_{S_0}(Y_0^3,N_0).
\end{aligned}
\right.
\end{equation}

对 \(c(X,Y,\mu^{2,Y})=u\) 两次应用\SGARefLink{sga3:iii:proposition:4-13}{4.13}, 得到
\[
 \begin{aligned}
 a_1\bigl(c(X^2,Y^2,f_1)\bigr)
 +b_1\bigl(c(X,Y,\mu_Y)\bigr)
 &=u\\
 &=a_2\bigl(c(X^2,Y^2,f_2)\bigr)
 +b_2\bigl(c(X,Y,\mu_Y)\bigr).
 \end{aligned}
\]
现在 \(c(X,Y,\mu_Y)=DY\), 且由于
\(f_1=(1,\mu)\) 及 \(f_2=(\mu,1)\), 采用显然的记号有
\[
 c(X^2,Y^2,f_1)=(0,DY),
 \qquad
 c(X^2,Y^2,f_2)=(DY,0).
\]
因此,
\[
 u=a_1((0,DY))+b_1(DY)
  =a_2((DY,0))+b_2(DY).
\]
首先注意, \(b_j\) 就是
\(\operatorname{Hom}_{S_0}((f_j^Y)_0,N_0)\),
即由 \((f_j^Y)_0\) 经函子性得到的态射. 因此上述恒等式成为
\[
 \begin{aligned}
 a_1((0,DY))
 &-\operatorname{Hom}((\mu,1),N_0)(DY)
 +\operatorname{Hom}((1,\mu),N_0)(DY)\\
 &\hspace{9em}-a_2((DY,0))=0.
 \end{aligned}
\]
中间两项就是 \(2\)-余边界公式中的
\(DY(xy,z)\) 和 \(DY(x,yz)\) 两项.
现在只需确定另外两项.

先计算映射 \(a_j\). 它由下列 \(\mathscr O_{Y_0^2}\)-模态射
沿 \((f_j^Y)_0\) 拉回得到:
\[
 P:\mathbf n_{Y_0/X_0}\otimes\mathscr O_{Y_0^2}
 \longrightarrow
 (\mathbf n_{Y_0/X_0}\oplus \mathbf n_{Y_0/X_0})
 \otimes\mathscr O_{Y_0^2}
\]
其中此态射由 \(Y_0\) 中的乘法诱导.
考虑向量丛 \(V=\mathbf V(\mathbf n_{Y_0/X_0})\).
由对偶性, \(P\) 给出
\[
 \mathbf V(P):
 V\times_{S_0}V\times_{S_0}Y_0\times_{S_0}Y_0
 \longrightarrow V\times_{S_0}Y_0\times_{S_0}Y_0,
\]
它在集合论意义下表示为
\[
 \mathbf V(P)(u,v,a,b)=(u+\operatorname{Ad}(a)v,ab,b).
\]
\footnote{编者注: 已将分量 \(a,b\) 换为 \(ab,b\),
以表明 \(\mathbf V(P)\) 是乘法态射
\(V_{Y_0}\times_{S_0}V_{Y_0}\to V_{Y_0}\) 拉回到 \(Y_0^2\) 所得.}
证明与 Lie 代数的相应事实, 即关于模 \(\omega^1_{Y_0/S_0}\) 的事实的证明完全相同.
首先注意, 由关于 \(Y_0\) 的函子性, \(V\) 在 \(S_0\) 上的向量丛范畴中带有群结构.
由已用于 Lie 代数的引理 (\SGARefLink{sga3:ii:lemma:3-10}{第 II 篇, 3.10}),
此结构与其 \(\mathscr O_S\)-模结构的底群结构一致. 其次,
\[
 \mathbf V\bigl(\mathbf n_{Y_0/X_0}
   \otimes_{\mathscr O_{S_0}}\mathscr O_{Y_0}\bigr)
 =\mathbf V(\mathscr N_{Y_0/X_0})
\]
自身也带有 \(S_0\)-群结构, 即 \(V\) 上的群结构与
\(Y_0\) 上的群结构所给出的半直积.
只需确定 \(Y_0\) 在 \(V\) 上的作用, 即得所求公式.

% Source page: Exp3 local p. 65; combined p. 175; printed pp. 145--146.
现在计算余下两项. 先考虑 \(a_1((0,DY))\).
它由下图计算, 其中 \(\mathbf n\) 表示 \(\mathbf n_{Y_0/X_0}\):

% Raster stand-in for Loop 1; replace by a native TeX diagram in Loop 2.
% Authority: Exp3 local p. 65; combined reader p. 175; printed p. 146.
\begin{center}
 \SGAThreeNativeInput[0.51\textwidth]
   {figures/exp3/diagram_057_exp3_p065_module_pullback.tex}
\end{center}

将这些模定义的向量丛视为 \(S_0\) 上的概形,
并取以任意概形为值的点. 若 \((u,x,y,z)\) 表示
\(\mathbf V(\mathscr J)\times Y_0^3\) 的一点, 则得到下图:

% Raster stand-in for Loop 1; replace by a native TeX diagram in Loop 2.
% Authority: Exp3 local p. 65; combined reader p. 175; printed p. 146.
\begin{center}
 \SGAThreeNativeInput[0.78\textwidth]
   {figures/exp3/diagram_058_exp3_p065_vector_bundle_points.tex}
\end{center}

因此
\[
 a_1((0,DY))(x,y,z)
 =\operatorname{Ad}(x)DY(y,z),
\]
这正是余边界的第一项. 同样,
\(a_2((DY,0))(x,y,z)=DY(x,y)\), 因此
\[
 \begin{aligned}
 0={}&\operatorname{Ad}(x)DY(y,z)-DY(xy,z)
      +DY(x,yz)-DY(x,y)\\
   ={}&(\partial^2DY)(x,y,z).
 \end{aligned}
\]
\footnote{编者注: 修改了符号, 使之与第 I 篇\SGARefLink{sga3:i:subsection:5-1}{5.1} 一致.}

% Source page: Exp3 local p. 66; combined p. 176; printed pp. 146--147.
\medskip
\noindent\SGAInlineRefTarget{sga3:iii:paragraph:4-20}\textbf{4.20. --- 证明 \SGARefLink{sga3:iii:lemma:4-17:item:iii}{4.17(iii)}.}
\footnote{编者注: 下面的论证在原版基础上作了扩充.}
比较 \(v=c(X,Y,\mu\circ(i',i'))\) 在下述两个图中算得的两个值:

% Raster stand-in for Loop 1; replace by native TeX diagrams in Loop 2.
% Authority: Exp3 local p. 66; combined reader p. 176; printed p. 147.
\begin{center}
 \SGAThreeNativeInput[0.78\textwidth]
   {figures/exp3/diagram_059_exp3_p066_comparison_diagrams.tex}
\end{center}

设 \(f=\mu\circ(i',i')\). 图 \((*)\) 给出
\begin{equation}
 v=DY'+f_0^*\bigl(c(X,Y,i')\bigr).
 \tag{1}\label{sga3:iii:paragraph:4-20:equation:1}
\end{equation}
现在 \(Y'_0=Y_0\), 且 \(f_0\) 为乘法
\(Y_0^2\to Y_0\); 故
\begin{equation}
 f_0^*\bigl(c(X,Y,i')\bigr)(x_0,y_0)
 =c(X,Y,i')(x_0y_0).
 \tag{2}\label{sga3:iii:paragraph:4-20:equation:2}
\end{equation}
令 \(c=c(X,Y,i')\). 在等同
\(N'_0\simeq N_0\oplus N_0\) 下, 元素
\[
 c(X\times_S X,Y\times_S Y,(i',i'))
\]
与对 \((c,c)\) 等同. 因此, 令 \(h=(i',i')\), 图 \((\dagger)\) 给出
\begin{equation}
 v=h_0^*(DY)+a_{\mu_0}(c,c).
 \tag{3}\label{sga3:iii:paragraph:4-20:equation:3}
\end{equation}
由于 \(h_0\) 是 \(Y_0^2\) 上的恒等映射, 有 \(h_0^*(DY)=DY\).
最后, 由前面对 \(a_{\mu_0}\) 的计算, 对每个 \(S'\to S\) 及每个
\(x_0,y_0\in Y_0(S'_0)\),
\begin{equation}
 a_{\mu_0}(c,c)(x_0,y_0)
 =c(x_0)+\operatorname{Ad}(x_0)\bigl(c(y_0)\bigr).
 \tag{4}\label{sga3:iii:paragraph:4-20:equation:4}
\end{equation}
从而
\begin{equation}
\begin{aligned}
 (DY'-DY)(x_0,y_0)
  ={}&\operatorname{Ad}(x_0)\bigl(c(X,Y,i')(y_0)\bigr)
      -c(X,Y,i')(x_0y_0)+c(X,Y,i')(x_0)\\
  ={}&\bigl(\partial^1c(X,Y,i')\bigr)(x_0,y_0).
\end{aligned}
\tag{5}\label{sga3:iii:paragraph:4-20:equation:5}
\end{equation}
由于 \(c(X,Y,i')=d(Y,Y')\)
(参见\SGARefLink{sga3:iii:proposition:4-8:item:ii}{4.8(ii)}), 这就证明了
\[
 DY'-DY=\partial^1d(Y,Y').
\]

\medskip
\noindent\SGAInlineRefTarget{sga3:iii:theorem:4-21}\textbf{定理 4.21.}
\textit{设 \(S\) 为概形, \(\mathscr I\) 和 \(\mathscr J\) 为 \(S\) 上的两个理想\footnote{编者注: 已删除 \(\mathscr I\) 幂零的假设, 因为它似乎是多余的.} 使得 \(\mathscr I\supset\mathscr J\) 且 \(\mathscr I\mathscr J=0\).
设 \(X\) 为在 \(S\) 上平滑的 \(S\)-群,
\(Y_{\mathscr J}\) 为 \(X_{\mathscr J}\) 的子 \(S_{\mathscr J}\)-群,
在 \(S_{\mathscr J}\) 上平坦且局部有限呈示.
考虑由下式定义的交换群值 \(S_0\)-函子 \(N_0\):}
\[
 \begin{aligned}
 \operatorname{Hom}_{S_0}(T,N_0)
 ={}&\operatorname{Hom}_{\mathscr O_T}
 \left(
  \mathbf n_{Y_0/X_0}\otimes_{\mathscr O_{S_0}}\mathscr O_T,
  \mathscr J\otimes_{\mathscr O_{S_0}}\mathscr O_T
 \right),\\[-0.2em]
 &\hspace{17em}T\in\operatorname{Ob}(\mathbf{Sch}/S_0),
 \end{aligned}
\]
\textit{\(Y_0\) 通过 \(X_0\) 的内自同构作用于它.}
\begin{enumerate}
\item[(i)] \textit{存在 \(X\) 的在 \(S\) 上平坦且约化为
\(Y_{\mathscr J}\) 的子 \(S\)-群的充分必要条件是下列两个条件成立:}
  \begin{enumerate}
  \item[\((i_1)\)] \textit{存在 \(X\) 的在 \(S\) 上平坦且提升 \(Y_{\mathscr J}\) 的子概形 \(Y\).
当 \(Y_0\) 仿射时, 此条件自动满足 (参见\SGARefLink{sga3:iii:proposition:4-6-5}{4.6.5}).}
  \item[\((i_2)\)] \textit{某个属于 \(H^2(Y_0,N_0)\) 的典范障碍消失.}
  \end{enumerate}
\item[\SGAInlineRefTarget{sga3:iii:theorem:4-21:item:ii}(ii)] \textit{若 (i) 中的条件成立, 则 \(X\) 的在 \(S\) 上平坦且约化为
\(Y_{\mathscr J}\) 的子 \(S\)-群 \(Y\) 所组成的集合,
是 \(Z^1(Y_0,N_0)\) 下的主齐性集.}\footnote{编者注: 对那些诱导 \(X(S_{\mathscr J})\) 的单位元的元素
\(x\in X(S)\) 的共轭作用取商后, 此集合是否是
\(H^1(Y_0,N_0)\) 下的主齐性集,
这一问题将在\SGARefLink{sga3:iii:paragraph:4-23}{4.23}--\SGARefLink{sga3:iii:proposition:4-36}{4.36} 讨论.}
\end{enumerate}

% Source page: Exp3 local p. 67; combined p. 177; printed pp. 147--148.
事实上, 条件 \((i_1)\) 是必要的. 假设它成立,
并设 \(Y\) 在 \(S\) 上平坦且提升 \(Y_{\mathscr J}\).
须找到同样在 \(S\) 上平坦且提升 \(Y_{\mathscr J}\) 的 \(Y'\),
使得 \(DY'=0\),\footnote{编者注: 已将原版的表达式 \(\partial DY'\) 修正为 \(DY'\).} 参见\SGARefLink{sga3:iii:lemma:4-17:item:i}{4.17(i)}.
由\SGARefLink{sga3:iii:lemma:4-17:item:iii}{4.17(iii)}, 这等价于找到元素
\(d(Y',Y)\in C^1(Y_0,N_0)\) 使得
\[
 DY=\partial^1d(Y',Y).
\]
\footnote{编者注: 见\SGARefLink{sga3:iii:lemma:4-17:item:iii}{4.17(iii)} 的编者注.}
设 \(c\in H^2(Y_0,N_0)\) 为 \(DY\) 所代表的类,
由\SGARefLink{sga3:iii:lemma:4-17:item:ii}{4.17(ii)}, 后者为余循环.
由\SGARefLink{sga3:iii:lemma:4-17:item:iii}{4.17(iii)}, 此类与 \(Y\) 的选择无关,
且其消失是存在满足上述方程的 \(d(Y',Y)\) 的充分必要条件.
这就证明了 (i).

若现在选定 \(Y\) 使 \(DY=0\), 则待解方程成为
\(\partial^1d(Y',Y)=0\), 这就证明了\SGARefLink{sga3:iii:theorem:4-21:item:ii}{(ii)}.

\medskip
\noindent\SGAInlineRefTarget{sga3:iii:remark:4-22}\textbf{注 4.22.}
沿用\SGARefLink{sga3:iii:theorem:4-21}{4.21} 的记号.
由\SGARefLink{sga3:iii:proposition:4-15}{4.15}, \(Y_0\) 在 \(X_0\) 中是局部完全交;
故 \(\mathscr N_{Y_0/X_0}\) 为有限局部自由 \(\mathscr O_{Y_0}\)-模, 因此
\[
 \mathbf n_{Y_0/X_0}
 =\varepsilon_0^*(\mathscr N_{Y_0/X_0})
\]
为有限局部自由 \(\mathscr O_{S_0}\)-模. 因此, 记
\[
 \mathbf n_{Y_0/X_0}^{\vee}
 =\operatorname{Hom}_{\mathscr O_{Y_0}}
  (\mathbf n_{Y_0/X_0},\mathscr O_{Y_0}),
\]
则有
\[
 \begin{aligned}
 &\operatorname{Hom}_{\mathscr O_T}
 \left(
  \mathbf n_{Y_0/X_0}\otimes_{\mathscr O_{S_0}}\mathscr O_T,
  \mathscr J\otimes_{\mathscr O_{S_0}}\mathscr O_T
 \right)\\
 &\qquad\simeq
 \mathbf n_{Y_0/X_0}^{\vee}
 \otimes_{\mathscr O_{S_0}}\mathscr J
 \otimes_{\mathscr O_{S_0}}\mathscr O_T
 \end{aligned}
\]
其中 \(T\to S_0\) 任意.\footnote{编者注: 扩充了上面的论证. 它表明下述同构不需要平坦性假设也成立.
另一方面, 从\SGARefLink{sga3:iii:paragraph:4-16}{4.16} 起,
讨论限制在 \(S\) 上平坦的 \(S\)-概形 \(f:T\to S\),
以保证障碍 \(c(X,Y,f)\) 所在的群
\[
 \operatorname{Hom}_{\mathscr O_{T_0}}
 \left(
  \mathbf n_{Y_0/X_0}\otimes_{\mathscr O_{S_0}}
   \mathscr O_{T_0},
  \mathscr J\mathscr O_T
 \right),
\]
与 \(N_0(T_0)\) 一致;
见\SGARefLink{sga3:iii:paragraph:4-16}{4.16} 末尾.}
因此, \(S_0\)-函子 \(N_0\) 同构于
\[
 \mathbf W\left(
  \mathbf n_{Y_0/X_0}^{\vee}
  \otimes_{\mathscr O_{S_0}}\mathscr J
 \right)
 \simeq
 \mathbf W\left(
  \operatorname{Hom}_{\mathscr O_{S_0}}
  (\mathbf n_{Y_0/X_0},\mathscr J)
 \right).
\]
从而存在同构\footnote{编者注: 采用第 I 篇\SGARefLink{sga3:i:subsection:5-3}{5.3} 的记号,
并假设 \(Y_0\) 在 \(S_0\) 上仿射.}
\[
 \begin{aligned}
 H^2(Y_0,N_0)
 &\simeq H^2\!\left(
   Y_0,\operatorname{Hom}_{\mathscr O_{S_0}}
   (\mathbf n_{Y_0/X_0},\mathscr J)\right)\\
 &\simeq H^2\!\left(
   Y_0,\mathbf n_{Y_0/X_0}^{\vee}
   \otimes_{\mathscr O_{S_0}}\mathscr J\right),\\[0.4em]
 Z^1(Y_0,N_0)
 &\simeq Z^1\!\left(
   Y_0,\operatorname{Hom}_{\mathscr O_{S_0}}
   (\mathbf n_{Y_0/X_0},\mathscr J)\right)\\
 &\simeq Z^1\!\left(
   Y_0,\mathbf n_{Y_0/X_0}^{\vee}
   \otimes_{\mathscr O_{S_0}}\mathscr J\right).
 \end{aligned}
\]

% Continue with 4.23 on Exp3 local p. 67 / combined p. 177 /
% printed p. 148.

% Translation coverage: entire assigned source component, 490 source lines.
% Formula audit: 223 expressions; math unchanged except translated text annotations.
% All source reference keys, labels, footnotes and native diagram calls retained.
